

%%%%%%%%%%%%%%%
\begin{frame}[fragile]{Le problème}

Représentation relationnelle  des entités\,: nuplets (entités),
et mécanisme clés primaire / clé étrangère (liens entre entités).

 
\vfill

Les applications objet (Java, C\#, Python, PHP, ...) représentent les entités
comme des \alert{objets} directement liés les uns aux autres.


\vfill

\alert{Problème}\,: comment \alert{convertir}  d'une représentation à l'autre\,?


\begin{blocgauche}
\begin{remblock}{Bases objets}
Une idée serait d'utiliser des bases \alert{objets} et
pas relationnelles. Elle n'a pas vraiment abouti.
\end{remblock}

\end{blocgauche}
\end{frame}


%%%%%%%%%%%%%%%
\begin{frame}[fragile]{Premier exemple de conversion}

Pour lire les contacts comme des objets d'une classe \texttt{Contact}.

\vfill

\begin{minted}[baselinestretch=.9,
               fontsize=\small,
               framesep=2mm]{python}
class Contact:
     def __init__(self,id,prenom,nom,email):
                self.id=id
                self.prenom=prenom
                self.nom=nom
                self.email=email

curseur.execute("select * from Contact")
for cdict in curseur.fetchall():
    # Conversion du dictionnaire en objet
    cobj = Contact(cdict["idContact"], cdict["prénom"], 
                cdict["nom"], cdict["email"])
    print(cobj.prenom, cobj.nom)
\end{minted}


\vfill

\begin{blocgauche}
Implique des instructions lourdes, répétitives,
sujettes à erreur.
\end{blocgauche}
\end{frame}


%%%%%%%%%%%%%%%
\begin{frame}[fragile]{Plus compliqué\,: la classe  \texttt{Message}}

Composition d'objets,  associations, méthodes.

\vfill

\vfill

\begin{minted}[baselinestretch=.9,
               fontsize=\small,
               framesep=2mm]{python}
class Message:
    # Emetteur: un objet 'Contact'
    emetteur = Contact(0, "", "", "")
    # Prédecesseur: peut ne pas exister
    predecesseur = None
    # Liste des destinataires = des objets 'Contacts'
    destinataires = []

    # La méthode d'envoi de message
    def envoi(self):
        for dest in self.destinaires:
           ...
\end{minted}

\vfill

Conversion depuis la base\,? Très pénible....
\end{frame}

%%%%%%%%%%%%%%%
\begin{frame}{Comment une application objet ``voit'' les données}

Comme un graphe\,!

\vfill

\vfill

\figSlideWithSize{instance-messagerie.pdf}{9cm}

\vfill

\begin{blocgauche}
Trop différent de la représentation relationnelle\,: il nous faut des outils.

\end{blocgauche}
\end{frame}


%%%%%%%%%%%%%%%
\begin{frame}{Le \textit{mapping} objet-relationnel (ORM)}

Un système ORM convertit \alert{automatiquement}, 
\alert{à la demande}, la base relationnelle sous forme d’un graphe d’objet

\vfill

\figSlideWithSize{mapping.pdf}{8cm}

\end{frame}

%%%%%%%%%%%%%%%
\begin{frame}[fragile]{Exemple avec le \textit{framework} Hibernate (Java)}

\begin{minted}[baselinestretch=.9,
               fontsize=\footnotesize,
               framesep=2mm]{java}
@Entity(table="Message")
public class Message {
  @Id
  private Integer id;
  
  @Column
  private String contenu;
  
  @ManyToOne
  private Contact emetteur;
  
  @OneToMany
  private Set<Contact> destinataires ;
}
\end{minted}


\vfill

\begin{blocgauche}
On indique au système ORM tout ce qui est nécessaire pour associer
objets et nuplets de la base.
\end{blocgauche}
\end{frame}


%%%%%%%%%%%%%%%
\begin{frame}[fragile]{Parcours d'un graphe d'objets}

Exemple (un peu simplifié) d'envoi des messages de Serge. 
\vfill
\begin{minted}[baselinestretch=.9,
               fontsize=\small,
               framesep=2mm]{java}
List<Message> resultat =
     session.execute("from Message as m "
                      + " where m.emetteur.prenom = 'Serge'");
                      
for (Message m : resultat) {
        for (Contact c : m.destinataires) {
                      message.envoi (c.email);
        }
}
\end{minted}

\vfill

Le système engendre et exécute les requêtes SQL. 

\vfill

\begin{blocgauche}
Tendance à 
produire \alert{beaucoup} de requêtes élémentaires là où \alert{une seule}
jointure
serait plus efficace.
\end{blocgauche}
\end{frame}

%%%%%%%%%%%%
\begin{frame}{À retenir}

Les ORMs sont maintenant très utilisés pour les développements d'envergure
\vfill

\begin{redbullet}
  \item Les accès à la base prennent la forme d'une navigation dans un graphe d'objets
  \item Le système engendre les requêtes SQL pour matérialiser le graphe
  \item Conversion entièrement automatique
  \item Gains de productivité, \alert{mais} utilisation sous-optimale de SQL
\end{redbullet}

\vfill

\begin{blocgauche}
Pour aller plus loin\,: la mini-application Django sur \url{http://sql.bdpedia.fr}
ou le cours complet \url{http://orm.bdpedia.fr}
\end{blocgauche}
\end{frame}

